\chapter{Introduction}

\section{Free-Space Optical Communication and Turbulence}

Free-space optical (FSO) communication offers fiber-like bandwidth without a
physical waveguide, but the atmosphere couples amplitude scintillation,
wavefront distortion, and beam wander into the received field. In this report
we focus on a passive optical receiver question: can a fixed diffractive neural
network improve detector-plane utility for a turbulent laser beam without any
adaptive feedback?

The turbulence models used throughout the campaign follow the standard
Kolmogorov and von Karman descriptions,
\begin{equation}
\Phi(\kappa) = 0.023\, r_0^{-5/3}\, \kappa^{-11/3},
\end{equation}
\begin{equation}
\Phi(\kappa) = 0.023\, r_0^{-5/3}\, (\kappa^2 + \kappa_0^2)^{-11/6},
\end{equation}
where $r_0$ is the Fried parameter, $\kappa$ is the spatial frequency, and
$\kappa_0$ is the outer-scale cutoff. The quantities $D/r_0$ and the Rytov
variance $\sigma_R^2$ will appear repeatedly because they provide a compact
regime description for weak, moderate, and strong turbulence.

\section{Phase-Only D2NN and Propagation Model}

The optical processor is modeled as a stack of phase-only diffractive layers.
Each layer applies
\begin{equation}
t(x,y) = \exp\!\bigl(j \phi(x,y)\bigr),
\end{equation}
and free-space propagation between layers is simulated with the band-limited
angular spectrum method (BL-ASM). In the sampled scalar model used in the
project, a phase-only D2NN block is a lossless linear operator up to numerical
precision, which is the starting point for the invariance arguments in
Chapter~4.

For detector-plane metrics the field is optionally relayed through a focal lens.
This distinction is central to the report:
\begin{itemize}
  \item output-plane metrics: $\CO$, wavefront RMS, phase residuals;
  \item focal-plane metrics: $\PIB$, Strehl, and received power $\RP@10\,\mu$m.
\end{itemize}
The campaign only converged once these two measurement planes were kept
strictly separate.

\section{Notation and Measurement Conventions}

The report reuses a small set of quantities throughout, so we fix them here.
Unless otherwise stated, all inner products and norms are taken on the sampled
receiver grid after the telescope aperture and before detector cropping.
\begin{align}
\CO(U_1,U_2)
&=
\frac{|\langle U_1,U_2\rangle|}{\|U_1\|_2\,\|U_2\|_2},\\
\PIB(r_b)
&=
\frac{\int_{|\rho|<r_b} I_f(\rho)\,d\rho}{\int I_f(\rho)\,d\rho},\\
\TP
&=
\frac{\int |U_{\mathrm{out}}|^2\,dA}{\int |U_{\mathrm{in}}|^2\,dA},\\
\RP(r_b)
&=
\int_{|\rho|<r_b} I_f(\rho)\,d\rho .
\end{align}
Here $U_{\mathrm{in}}$ and $U_{\mathrm{out}}$ are the sampled complex fields
before and after the diffractive block, while $I_f$ is the focal-plane
irradiance after the detector lens. The shorthand $\RP@10\,\mu$m used in later
tables means the bucket power with $r_b=10\,\mu$m.

\section{Research Question}

The campaign began with an optimistic hypothesis: if turbulence is a structured
optical distortion, perhaps a static diffractive network can be trained as a
wavefront corrector. The experiments instead led to a narrower and more
defensible conclusion.

\begin{keyinsight}
Static phase-only D2NNs do not improve field-level similarity under the sampled
unitary model, but they can reshape intensity observables at the detector. The
entire report is the chronology of discovering where that distinction matters
and where it breaks naive interpretations.
\end{keyinsight}

This distinction also motivates the Received Power objective used later in the
campaign. A detector does not observe $\CO$ directly; it observes power falling
into a finite aperture or bucket. Once that detector-centric formulation is
made explicit, the design space changes substantially.

\section{Report Roadmap}

Chapters~2--3 cover the first spatial D2NN beam-reducer experiments and the
FD2NN optics crisis around the $f=1$\,mm design. Chapters~4--5 establish the
static-limit theorem and diagnose why the original PIB interpretation was
misleading. Chapters~6--9 rebuild the pipeline around clean focal-plane metrics
and progressively refine the loss from normalized sharpness objectives to raw
received power. Chapter~10 studies distance-dependent regimes, while
Chapters~11--12 revisit a 4f FD2NN architecture under the corrected detector
objective and assess its limits conservatively. Chapter~13 summarizes the final
lessons and future directions.
